\section{VERA}
\label{sec:method}

\subsection{Overview}

VERA treats open-ended functional model induction as evidence production under
two linked risks. The agent must act on an unfamiliar application to discover
functionality, yet a proposed function or an executor-reported completion may
not establish that the expected application-level outcome occurred. At the
same time, the interactions used to obtain that evidence can alter the
environment. The method therefore governs both how functional claims become
persistent knowledge and how the actions used to test them are represented.

This view yields three design principles. First, \emph{claims precede
evidence}: VERA states a testable functional hypothesis and freezes its
expected observable outcome before execution. Second, \emph{knowledge requires
outcome evidence}: executor status remains a diagnostic signal, while admission
depends on whether the resulting observations support the frozen expectation.
Third, \emph{evidence acquisition is accountable interaction}: VERA assesses a
selected semantic browser action in its current GUI context before execution
and links the resulting risk record to the action, claim, and outcome evidence.

VERA operationalizes these principles in a persistent exploration cycle.
Candidate generation proposes claims; pre-execution risk assessment records
potential consequences; browser interaction produces observable evidence; and
post-execution verification governs admission to the functional model. A
persistent frontier, dependency-aware selection, and replay preserve and
recover unresolved verification opportunities so that evidence production can
continue across locations. The candidate generator, browser executor, risk
judge, and outcome verifier are replaceable components. VERA's method is the
protocol that connects their outputs into a traceable functional model, rather
than a new vision-language model or browser automation backend.

\subsection{Testable Function Hypothesis Proposal}

To evaluate evidence rather than reinterpret it after the fact, VERA first
turns a possible function into an explicit, testable claim. At a semantic
location $\ell$, the proposal component conditions on the current GUI
observation and the accumulated model to specify a high-level function together
with its expected observable outcome. This representation targets
application-level functionality---for example, adding an item or creating a
project---rather than treating every visible widget as a separate function.

The expected outcome defines in advance what subsequent evidence must support.
It may describe equivalent visible manifestations, but cannot rely solely on
hidden backend state. VERA stores the function and expectation together before
browser interaction and does not revise the expectation in response to the
observed result. The resulting hypothesis is registered as unresolved until
execution evidence is judged. Proposal generation therefore creates a claim to
test, not knowledge to admit.

VERA keeps the outcome claim distinct from metadata describing where and when
it has been observed to apply. Location constraints associate the claim with
its semantic context. When testing a function relies on an earlier interaction,
VERA also records that relation as an observed direct dependency. Such metadata
helps select currently executable hypotheses and preserve local interaction
context, but does not assert complete business preconditions or causal rules.

\subsection{Pre-Execution Risk Awareness}

Evidence acquisition is an interaction with the environment, not a neutral
internal computation. VERA therefore assesses the selected semantic browser
action before it is executed. Because the consequence of an action depends on
the visible application state, the assessed object is the pair of the current
GUI observation and the selected action, rather than a context-free function
name. The same action label can consequently receive different judgments in
different states.

For each selected action, VERA captures the current screenshot and combines it
with the action label and a versioned risk taxonomy. An independent
vision-language judge returns a structured record containing a binary
\texttt{potential\_risk} decision, one primary \texttt{risk\_type} when risk is
identified, and brief interface-grounded evidence. The taxonomy supplies a
shared vocabulary of category definitions and examples; it is not a prediction
model on its own.

VERA attaches this record to the assessed action and, for an ordinary
verification attempt, to the functional claim and subsequent outcome evidence.
Replay actions receive separate records at the same semantic-action
granularity. These records expose a pre-execution signal that a downstream
policy could use for blocking or human confirmation. In the evaluated system,
however, the signal is recorded without changing the selected action; if risk
inference fails, execution continues and the error remains in the audit trail.

\subsection{Execution and Evidence Collection}

Testing a functional claim requires evidence about the application-level
outcome, not merely evidence that an interaction command was issued. The
browser executor receives the selected semantic action and attempts to realize
it. VERA stores the pre-execution observation, the concrete interaction trace,
the executor-reported status, the post-execution observation, URLs, and any
available structured state changes or evidence references. Together these
fields form the evidence packet in Section~\ref{sec:problem-formulation}.

This packet deliberately preserves both executor status and observable outcome
evidence. A successful executor status can indicate that a click or input
sequence completed even when the expected application-level change did not
occur. Conversely, a reported failure may reflect grounding, timing, or
environment errors rather than the absence of the proposed function. VERA uses
the status diagnostically and defers knowledge admission until the evidence is
compared with the frozen outcome claim.

\subsection{Evidence-Grounded Model Update}

Knowledge admission is a separate decision from both proposal and execution.
The outcome verifier receives the frozen functional claim and its evidence
packet, compares the expected outcome with the observed post-interaction state,
and records a \textsc{Supported}, \textsc{Unsupported}, or
\textsc{Unresolved} judgment. The current implementation uses a vision-language
model, but the admission protocol does not depend on a particular verifier
architecture.

VERA admits a hypothesis only when the required evidence is complete and the
evidence judgment is \textsc{Supported}, following
Equation~\ref{eq:admission}. The admitted record links the functional claim and
frozen expectation to the execution trace, before--after observations, evidence
judgment, and admission decision. Evidence that is missing, inconclusive, or
inconsistent with the expectation does not enter the downstream-usable
knowledge set. The associated claim and attempt nevertheless remain in the
model with their judgment and diagnostic reason. This separation prevents
unresolved execution problems from becoming functional facts while preserving
them as future verification opportunities.

\subsection{Persistent Evidence Production}

Conservative admission is useful only if exploration can continue producing
evidence for claims that are not resolved immediately. An interaction may move
the browser away from a claim's semantic location, an execution may fail, or a
run may exhaust its current budget. VERA retains such hypotheses in a
persistent frontier rather than equating interruption with rejection or
discarding the verification opportunity.

The selection procedure prioritizes hypotheses whose observed direct
dependencies are satisfied in the current context. When unresolved hypotheses
remain at a previously visited location, VERA may replay an observed path to
that location and resume evidence production. Replay restores a previously
observed context; it does not relax the evidence judgment or admission rule.

Replay is therefore a mechanism for recovering verification opportunities, not
an uncounted exploration action. VERA records replay interactions and their
pre-execution risk annotations separately from ordinary candidate attempts.
This separation permits an evaluation to hold the ordinary attempt budget
fixed while reporting the additional GUI cost and any supported functionality
obtained after replay.

Persistent frontier, dependency-aware selection, and replay implement the
evidence-production side of the method. They do not change the admission rule:
regardless of how a candidate is reached, only an evidence-supported outcome
can become admitted functional knowledge.

\subsection{Traceable Functional Model}

The resulting model preserves more than a list of discovered functions. For
each claim, it links the frozen expectation, selected semantic action,
pre-execution risk record, interaction trace, before--after evidence, evidence
judgment, and admission decision. Supported, unsupported, and unresolved
claims remain distinguishable, together with diagnostic reasons for
non-admission. This trace is the common representation through which VERA
accounts for both how evidence was acquired and why a claim did or did not
become reusable knowledge.

Downstream consumers can use a conservative view containing only admitted
functions, their semantic locations, observed direct dependencies, and
supported outcomes, while the complete trace remains available for inspection
or further verification. This interface does not require a particular planner
and does not treat the induced model as a complete recovery of hidden business
state or causal dynamics.
